\chapter{1967:Manfred Eigen and Ronald G.W. Norrish and George Porter}

\label{chap:chemistry1967}

The Nobel Prize in Chemistry for the year 1967 was awarded jointly to Manfred Eigen and Ronald G.W. Norrish and George Porter.

\textbf{Motivation:}

for their studies of extremely fast chemical reactions, effected by disturbing the equilibrium by means of very short pulses of energy

\section{Manfred Eigen}

\subsection*{Early Life and Education}

Manfred Eigen was born on 1927-05-09 in Bochum, Germany.

\subsection*{Career and Major Research}

Eigen studied physics and chemistry at the University of Göttingen, where he received his doctorate in 1951. He joined the Max Planck Institute for Physical Chemistry in Göttingen, later renamed the Max Planck Institute for Biophysical Chemistry, and became its director. He developed relaxation methods, including temperature-jump and pressure-jump techniques, to measure the rates of very fast reactions in solution.

\subsection*{Nobel Prize-winning Work}

The work for which Manfred Eigen was awarded the Nobel Prize in Chemistry in 1967 was described by the Nobel Committee as: "for their studies of extremely fast chemical reactions, effected by disturbing the equilibrium by means of very short pulses of energy".

Eigen disturbed a reaction equilibrium with a rapid change of temperature, pressure, or electric field and then followed the return to equilibrium, measuring rate constants for processes such as proton transfer that occur in microseconds.

\subsection*{Significance and Impact}

Eigen's relaxation methods opened the study of elementary chemical reactions in solution, including acid-base and enzyme-catalyzed reactions, to direct measurement.

\subsection*{Later Career and Honors}

Eigen later turned to the origin of life and to self-organization in biological systems. He died on 2019-02-06 in Göttingen, Germany.

\subsection*{Legacy}

Relaxation techniques remain standard tools in chemical kinetics, and Eigen's later work influenced research on the molecular basis of early evolution.

\subsection*{Selected Publications}

"Immeasurably Fast Reactions" (Nobel Lecture, 1967)

\clearpage

\section{Ronald G.W. Norrish}

\subsection*{Early Life and Education}

Ronald G.W. Norrish was born on 1897-11-09 in Cambridge, England.

\subsection*{Career and Major Research}

Norrish studied at the University of Cambridge, where he spent his academic career and became professor of physical chemistry. With George Porter he developed flash photolysis, in which an intense flash of light generates short-lived chemical species that are then analyzed by a second flash.

\subsection*{Nobel Prize-winning Work}

The work for which Ronald G.W. Norrish was awarded the Nobel Prize in Chemistry in 1967 was described by the Nobel Committee as: "for their studies of extremely fast chemical reactions, effected by disturbing the equilibrium by means of very short pulses of energy".

Flash photolysis made it possible to observe free radicals and other transient species directly and to measure the rates of their reactions.

\subsection*{Significance and Impact}

Flash photolysis transformed the study of photochemistry and provided the first direct observations of many reactive intermediates in the gas phase.

\subsection*{Later Career and Honors}

Norrish remained at Cambridge until his retirement in 1965. He died on 1978-06-04 in Cambridge, England.

\subsection*{Legacy}

The Norrish type I and type II photochemical reactions of carbonyl compounds remain standard topics in photochemistry.

\subsection*{Selected Publications}

"Chemical Reactions Produced by Very High Light Intensities" (Nature, 1949)

\clearpage

\section{George Porter}

\subsection*{Early Life and Education}

George Porter was born on 1920-12-06 in Stainforth, England.

\subsection*{Career and Major Research}

Porter studied at the University of Leeds and then at the University of Cambridge, where he received his doctorate under Norrish. He was professor at the University of Sheffield and later director of the Royal Institution in London. He extended flash photolysis to study very fast reactions in gases and liquids.

\subsection*{Nobel Prize-winning Work}

The work for which George Porter was awarded the Nobel Prize in Chemistry in 1967 was described by the Nobel Committee as: "for their studies of extremely fast chemical reactions, effected by disturbing the equilibrium by means of very short pulses of energy".

Porter's refinements of flash photolysis allowed the spectra and kinetics of short-lived molecules to be recorded, and he applied the method to fields from photochemistry to photosynthesis.

\subsection*{Significance and Impact}

Flash photolysis became a cornerstone of photochemistry and photobiology and led directly to later laser-based techniques for studying fast processes.

\subsection*{Later Career and Honors}

Porter became director of the Royal Institution and president of the Royal Society and was created Baron Porter of Luddenham. He died on 2002-08-31 in Canterbury, England.

\subsection*{Legacy}

Porter's work established the experimental study of fast photochemical processes and influenced fields from solar energy research to atmospheric chemistry.

\subsection*{Selected Publications}

"Chemical Reactions Produced by Very High Light Intensities" (Nature, 1949)

\clearpage

\section*{Chapter Summary}

The 1967 prize recognized Eigen, Norrish, and Porter for their studies of extremely fast chemical reactions. Eigen's relaxation techniques and Norrish and Porter's flash photolysis made the kinetics of elementary processes measurable.

